%------------------------------------------------------------------------------%

\usepackage{integracao_e_series}
\usepackage{systeme}

\title{Integração por Frações Parciais}
\subtitle{Caso 2 -- Fatores reais repetidos}

%------------------------------------------------------------------------------%
\begin{document}

\addtitle

%------------------------------------------------------------------------------%
\section{Preliminares}

%------------------------------------------------------------------------------%
\begin{frame}{Integral Relevante}
  \begin{align*}
    \onslide<+->{\int \frac{1}{x^2+a^2}\, dx}
    \onslide<+->{& = \int \frac{1}{a^2\left(\left(\frac{x}{a}\right)^2 + 1\right)} \,dx \\[2mm]}
    \onslide<+->{& = \frac{1}{a} \int \frac{1}{\left(\frac{x}{a}\right)^2 + 1} \,\frac{dx}{a}}
    \onslide<+->{& u = \frac{x}{a} \qquad du = \frac{dx}{a}  \\[2mm]}
    \onslide<+->{& = \frac{1}{a} \int \frac{1}{u^2+1}\, du  \\[2mm]}
    \onslide<+->{& = \frac{1}{a} \arctan(u) + C    \\[2mm]}
    \onslide<+->{& = \frac{1}{a} \arctan\left(\frac{x}{a}\right) + C}
  \end{align*}
\end{frame}

%------------------------------------------------------------------------------%
\section{Integração por Frações Parciais}

%------------------------------------------------------------------------------%
\begin{frame}{Caso 2}
  $Q(x)$ é o produto de $m$ fatores lineares e alguns repetidos

  \pause
  Suponha que o primeiro fator $(a_1 x + b_1)$ seja repetido $p$ vezes
  \[
    Q(x)
    = \left(a_1 x + b_1\right)^p
      \left(a_2 x + b_2\right) \cdots
      \left(a_q x + b_q\right)
  \]

  \pause
  Expansão
  \[
    \frac{P(x)}{Q(x)}
    = \frac{A_1}{(a_1x+b_1)}
    + \frac{A_2}{(a_1x+b_1)^2}
    + \cdots
    + \frac{A_p}{(a_1x+b_1)^p}
    + \cdots
    + \frac{A_q}{a_q x + b_q}
  \]
\end{frame}

\begin{frame}{Exemplo}
  \[
    \frac{x^3-x-2}{x^3(x-1)}
    = \frac{A}{x}
    + \frac{B}{x^2}
    + \frac{C}{x^3}
    + \frac{D}{x-1}
  \]

  \pause
  \vspace{\baselineskip}
  \[
    \int \frac{x^3-x-2}{x^3(x-1)} \,dx
    = A \int \frac{1}{x} \,dx
    + B \int \frac{1}{x^2} \,dx
    + C \int \frac{1}{x^3} \,dx
    + D \int \frac{1}{x-1} \,dx
  \]
\end{frame}

%------------------------------------------------------------------------------%
%------------------------------------------------------------------------------%
\addtocounter{exemplo}{1}
\section{Exemplo \theexemplo}

%------------------------------------------------------------------------------%
\begin{frame}{Exemplo \theexemplo}
  Calcule \(\int \frac{4x}{x^3-x^2-x+1}dx\)

  \pause
  O grau do numerador $P$ é menor que o grau do denominador $Q$

  \pause
  Raízes de $Q$: $x=1$ com multiplicidade 2 e $x=-1$

  \pause
  Fatorando $Q$
  \[
    Q(x)
    = x^3 - x^2 - x - 2
    = (x-1)^2(x+1)
  \]
  \pause
  Caso 2: Fatores lineares com repetição
\end{frame}

%------------------------------------------------------------------------------%
\begin{frame}{Exemplo \theexemplo}
  \onslide<+->{\[
    \frac{4x}{x^3-x^2-x-2}
    = \frac{A}{x-1}
    + \frac{B}{(x-1)^2}
    + \frac{C}{x+1}
  \]}
  \onslide<+->{Multiplicando os dois lados por \(x^3-x^2-x-2=(x-1)^2(x+1)\)}
  \begin{align*}
    \onslide<+->{4x
    & = A(x-1)(x+1)+B(x+1)+C(x-1)^2 \\}
    \onslide<+->{& = (A+C)x^2+(B-2C)x+(-A+B+C)}
  \end{align*}
  \onslide<+->{\systeme[ABC]{
     A    + C = 0,
        B -2C = 4,
    -A +B + C = 0
  }}
  \onslide<+->{\[
    A =  1 \qquad
    B =  2 \qquad
    C = -1
  \]}
\end{frame}

%------------------------------------------------------------------------------%
\begin{frame}{Exemplo \theexemplo}
  \begin{align*}
    \onslide<+->{\int \frac{4x}{x^3-x^2-x-2} \,dx
    & = \int \frac{1}{x-1}+ \frac{2}{(x-1)^2}-\frac{1}{x+1} \, dx \\[2mm]}
    \onslide<+->{& = \int \frac{1}{x-1} \,dx
      + \int \frac{2}{(x-1)^2} \,dx
      - \int \frac{1}{x+1} \,dx \\[2mm]}
    \onslide<+->{& = \ln \abs{x-1}- \frac{2}{x-1}- \ln\abs{x+1}+C \\[2mm]}
    \onslide<+->{& = \ln \abs{\frac{x-1}{x+1}} - \frac{2}{x-1}+C}
  \end{align*}
\end{frame}

%------------------------------------------------------------------------------%
\section{Lista Mínima}

%------------------------------------------------------------------------------%
\begin{frame}{Lista Mínima}
  Estudar a Seção 4.5 da Apostila

  Exercícios:

  \vfill
  Atenção: A prova é baseada no livro, não nas apresentações
\end{frame}

%------------------------------------------------------------------------------%
\end{document}
%------------------------------------------------------------------------------%